\section{Discussion}

Whole-cochlea inflammatory enrichment coexists with opposing responses in stromal, vascular and immune populations. TNF--NF-\(\kappa\)B enrichment favored noise exposure in fibrocyte/\allowbreak{}mesenchymal populations from both dissected regions. Endothelial and immune/\allowbreak{}myeloid populations showed the opposite direction, with neuronal and epithelial/\allowbreak{}supporting-associated populations adding regional variation. Animal omissions preserved these directions; fibrocyte FDR significance varied, while endothelial and myeloid enrichment remained significant.

Milon and colleagues established cell-type-specific transcriptional responses to acoustic trauma~\citep{Milon2021}. The fibrocyte response complements reporter evidence of noise-associated NF-\(\kappa\)B activity in spiral-ligament fibrocytes~\citep{Adams2009,Herranen2018}. The lateral-wall \mbox{\textit{Ccl5}} increase adds a chemokine-associated component to this response. Previous experiments detected cytokine production in spiral-ligament fibrocytes and chemokine and adhesion-molecule expression during leukocyte recruitment~\citep{Fujioka2006,Tornabene2006}. Together, these observations connect stromal inflammatory transcription with genes associated with local immune communication.

Myeloid TNF--NF-\(\kappa\)B enrichment was negative at 24 hours, within a temporally and anatomically heterogeneous immune response. Acoustic injury recruits cochlear mononuclear phagocytes and induces antigen-presentation features in basilar-membrane phagocytes~\citep{Hirose2005,Yang2015}. Resident macrophage responses have been linked to hair-cell injury and TLR4 signaling~\citep{Pan2024}. Fate mapping distinguishes the spatial and temporal distributions of resident and recruited macrophages after trauma~\citep{Murali2025}. Selected innate-immune genes increased in specific myeloid populations in the original study~\citep{Milon2021}. Their induction can coexist with negative enrichment of the broader TNF--NF-\(\kappa\)B set. Cell-state changes and subtype composition may both contribute to this program-specific myeloid response.

Proteasome-subunit genes comprised 28 of 48 MAPK-family leading-edge members in lateral-wall endothelial cells, alongside \mbox{\textit{Map2k2}}, \mbox{\textit{Mapk3}} and \mbox{\textit{Dusp2}}. This composition links the annotated signaling program to a substantial protein-quality-control transcriptional response. Whole-cochlea protein-folding and unfolded-protein-response enrichment fits the proteotoxic response described by Jongkamonwiwat and colleagues~\citep{Jongkamonwiwat2020}. Independent experiments have measured time-dependent MAPK phosphorylation~\citep{Maeda2013} and delayed cochlear oxidative damage after acoustic trauma~\citep{Yamashita2004}. NRF2 loss and pharmacological activation also link antioxidant defense to acoustic vulnerability in mice~\citep{Honkura2016}.

Endothelial stress responses also connect to the cellular architecture of the cochlear barrier. Interactions among cochlear endothelial cells, pericytes and perivascular macrophage-like melanocytes regulate strial permeability~\citep{Neng2013}. Noise exposure can alter perivascular contacts and endothelial junctions in this compartment~\citep{Zhang2013}.

Expression localization adds cell-type context to condition-associated changes across cochlear populations. Endothelial and fibrocyte/\allowbreak{}mesenchymal TNF--NF-\(\kappa\)B localization in GSE252285 coexisted with opposite enrichment directions in GSE168041. \mbox{\textit{Apold1}} expression spanned endothelial cells and fibrocyte-associated cluster 19, retaining an endothelial maximum in mean expression across broad cell types. Reassigning cluster 19 to fibrocytes preserved the eleven recurrent localization patterns and all matched-null significance classifications. \mbox{\textit{Apold1}} and \mbox{\textit{Errfi1}} thus illustrate recurrent localization despite nonsignificant condition effects in the replicated profiles.

The datasets sample distinct exposure windows and anatomical settings in the cochlea. GSE160639 sampled whole cochleae immediately after a 30-minute exposure, whereas GSE168041 profiled dissected compartments 24 hours after a two-hour exposure~\citep{Jongkamonwiwat2020,Milon2021}. GSE252285 offers descriptive localization from one pooled library per condition, collected six hours after the final exposure~\citep{GSE252285}. Its UCell-guided comparisons remain exploratory because the six-test FDR family excludes the preceding program--cell-type screen. Documented strain-dependent fibrocyte NF-\(\kappa\)B responses also inform comparisons across these experimental settings~\citep{Herranen2018}.

Fibrocyte-associated TNF--NF-\(\kappa\)B and chemokine transcription form a stromal component of the cochlear response to acoustic trauma. Endothelial and immune/\allowbreak{}myeloid populations show opposing enrichment of the same inflammatory gene set. Proteasome-rich endothelial MAPK enrichment and whole-cochlea protein-folding responses place protein quality control alongside inflammation in the cochlear stress response.
